\begin{abstract}
Bipolar Disorder (BD) is routinely mistaken for Major Depressive Disorder in clinical settings, largely because patients present during depressive crashes that look nearly identical across both conditions. This misdiagnosis is not merely inconvenient it can be dangerous, as antidepressants prescribed to bipolar patients can precipitate severe manic episodes. In this work, we explore whether continuous wrist-worn actigraphy data, captured over multi-day monitoring periods, carries enough signal to make this distinction automatically. Using the open-source Depresjon dataset, we develop a hybrid 1D Convolutional Neural Network and Long Short-Term Memory (1D-CNN-LSTM) architecture that operates directly on raw minute-level activity counts, without any hand-crafted feature engineering. Our two-stage experimental design first validates the model on a healthy-versus-depressed classification task, then tackles the harder and clinically meaningful challenge of separating bipolar from unipolar depressive episodes. We show that deep sequence models can capture both the short-term kinetic texture and the longer circadian rhythm disruptions that distinguish these two conditions.

\keywords{bipolar disorder, unipolar depression, actigraphy, CNN-LSTM, sequence modeling, wearable sensors, differential diagnosis}
\end{abstract}
